\chapter{A request, end to end}\label{ch:request-end-to-end}

\whatyoufind{One memory request followed from the core that issues it to the core that
retires it, drawn on the system as shipped and on a time axis; the Logger stamp at every hop;
and, from those stamps, the exact definition of every column of a \code{LatencyReport}. This is
the chapter to have open the first time you read a report: a total tells you a core was slow,
the columns tell you which shared resource made it slow, and this is what each column is.}

\section{The journey}

A core issues a \code{Message}. Its L1 admits it, looks the line up and, on a miss, puts a
\code{GetS} (read) or \code{GetM} (write) on the request lane of \code{bus[0]}, where the arbiter
elects it. Every other L1 snoops it; the LLC admits it into its queue. If the LLC has the line
and no L1 owns it, the data is emitted when the array port is granted; if an L1 owns it, that
L1 answers on the response lane instead; if the LLC does not have it, a read leaves on
\code{bus[1]}, DRAM serves it, the fill comes back to the front of the LLC's queue and the
response is emitted from it. The response travels the response lane to the requesting L1, which
fills and hands the data to the core.

\figpage{system_end_to_end.pdf}{The system as shipped with one LLC \emph{miss} from core 0
routed through it. Each hop carries the colour of the \code{LatencyReport} column that measures
it (amber: waiting; purple: a bus transfer; teal: being served) and the number of the Logger
stamp that bounds it, \ding{172}--\ding{181} being the report's nine latency columns in order; the
side panel lists the stamp pair behind every column.}{fig:system-miss}

\figpage{system_end_to_end_hit.pdf}{The same system with an LLC \emph{hit} from core 0. The route
stops at the LLC's array port, where the data is emitted at the grant; the memory bus and DRAM
are not on the path, so \col{L2-DRAM Bus} and \col{DRAM} are 0.}{fig:system-hit}

\Cref{fig:system-miss,fig:system-hit} are the same drawing with the two routes. Three things
on them are worth saying in words. The other cores' L1s see every request on the request lane,
and when one of them owns the line it answers on the response lane in the LLC's place, its
stamps then standing in for the LLC's. The service lane carries the LLC's back-invalidations on
an eviction and is on neither request's path. And a hit's \col{L2 Access} ends at the port
\emph{grant}: the data leaves then, and the read's \code{A\_LLC} cycles on the port are paid by
whoever waits for the port next --- their \col{L2 Access}, or a fill's write.

\section{The stamps}

At each hop the component stamps a Logger checkpoint: an \stamp{ENTER} when it takes custody of
the request, a \stamp{SERVICE} when it begins the actual work, an \stamp{EXIT} when it hands the
request on. Every stamp names the component's \emph{role} (\code{CPU}, \code{L1},
\code{REQ\_BUS}, \code{RESP\_BUS}, \code{LLC}, \code{MEM\_BUS}, \code{DRAM}) and is taken in the
originating core's clock, so a request's timeline is one ordered list of self-describing events
and the columns are differences of \emph{named} stamps, never inferences from position
(\cref{ch:monitoring} has the mechanism). \Cref{tab:stamps} lists who stamps what.

\begin{table}[htbp]\centering\small
\caption{Where each Logger stamp is taken.}\label{tab:stamps}
\begin{tabularx}{\textwidth}{@{}>{\ttfamily}p{0.3\textwidth}Y@{}}
\toprule \normalfont stamp & taken by \\ \midrule
CPU.ENTER & \code{CPU::addRequest}: the request is issued (and bound to its core) \\
L1.ENTER & \code{BaseController::processLogic} when an L1 admits a request from below \\
REQ\_BUS.EXIT & \code{BusController::broadcast} on \code{bus[0]}: the request is transmitted \\
LLC.ENTER & \code{BaseController::processLogic} when the LLC admits the bus request \\
MEM\_BUS.EXIT & \code{BusController} on \code{bus[1]}, once in each direction \\
DRAM.ENTER, DRAM.EXIT & \code{MainMemoryController}: request arrives, data leaves after \key{m\_memory\_latency}. MCsim stamps neither \\
LLC.SERVICE, LLC.EXIT & \code{hitAction} / \code{removePendingAndRespond}: port granted; response emitted \\
L1.SERVICE, L1.EXIT (before RESP\_BUS.EXIT) & \code{performWriteBack} when an \emph{owning} L1 supplies another core's request; the responder is then that L1, not the LLC \\
RESP\_BUS.EXIT & \code{BusController::send} on \code{bus[0]}: the response is transmitted \\
CPU.EXIT & \code{CPU::checkReceiveBuffer}: delivered; triggers the report row \\
\bottomrule
\end{tabularx}
\end{table}

\figpage{request_end_to_end.pdf}{The same two routes on a time axis: an LLC hit and an LLC
miss, with the nine \code{LatencyReport} columns bracketed between the Logger stamps that
bound them. Ticks are in order, not to scale. Below the miss axis, the array write of the fill
runs alongside the response and lands in no column.}{fig:request-timeline}

\section{The columns}

\begin{table}[htbp]\centering\small
\caption{The \code{LatencyReport} columns, in report order, as differences of stamps. The
nine latency columns tile exactly to \col{Total}; a stamp a request never reaches makes its
column 0.}\label{tab:columns}
\begin{tabularx}{\textwidth}{@{}lY@{}}
\toprule column & definition \\ \midrule
\col{CPU Latency} & \stamp{CPU.ENTER} $-$ \col{Ready Cycle}: the wait inside the core before the request could issue (the out-of-order window). Lies \emph{before} Total \\
\col{L1 Stall} & \stamp{L1.ENTER} $-$ \stamp{CPU.ENTER} \\
\col{Request Bus} & \stamp{REQ\_BUS.EXIT} $-$ \stamp{L1.ENTER}: the L1's miss handling and its wait for the request lane \\
\col{L2 Stall} & \stamp{LLC.ENTER} $-$ \stamp{REQ\_BUS.EXIT}: the wait in the LLC's queue \\
\col{L2 Access} & hit: \stamp{LLC.SERVICE} $-$ \stamp{LLC.ENTER}, the wait for the array port; miss: \stamp{LLC.SERVICE} $-$ \stamp{MEM\_BUS.EXIT}$_2$, the fill's turn at the front of the queue, typically one cycle \\
\col{Response Bus} & \stamp{RESP\_BUS.EXIT} $-$ \stamp{LLC.EXIT}: the clean response-lane leg \\
\col{L2-DRAM Bus} & miss: (\stamp{DRAM.ENTER} $-$ \stamp{LLC.ENTER}) $+$ (\stamp{MEM\_BUS.EXIT}$_2$ $-$ \stamp{DRAM.EXIT}): one column over two hops, out and back \\
\col{DRAM} & miss: \stamp{DRAM.EXIT} $-$ \stamp{DRAM.ENTER} \\
\col{L1 Access} & \stamp{CPU.EXIT} $-$ $\max$(\stamp{L1.ENTER}, \stamp{RESP\_BUS.EXIT}): an L1 hit's own service, or the fill and hand-off after the response-lane grant \\
\col{Total} & \stamp{CPU.EXIT} $-$ \stamp{CPU.ENTER} $=$ the sum of the eight stage columns above (CPU Latency excluded) \\
\col{Effective} & \stamp{CPU.EXIT} $-$ $\max$(previous request's finish, \stamp{CPU.ENTER}): the part of Total not overlapped with the previous request; its mean is the per-core average, its sum the true added execution time \\
\col{Oldest} & \stamp{CPU.EXIT} $-$ the cycle this request became the oldest outstanding request of its core; 0 if it never was. At OoO$=1$ it equals Total \\
\bottomrule
\end{tabularx}
\end{table}

Three things in \cref{tab:columns} are easy to misread from the names alone.

\paragraph{L2 Access means two different waits.} On a hit it ends at the array-port grant, not
after \code{A\_LLC} cycles: the data is emitted at the grant, and the port's busy time is paid by
whoever waits for the port next. On a miss it is the fill's turn at the LLC: each cycle one
message moves from the memory-side RX into the \emph{front} of the processing queue, a data
message is always ready, and every ready message is handled that same cycle --- so typically a
single cycle, plus any backlog in that one FIFO, and never a wait behind queued demand
requests. The response is emitted the moment the fill is processed, carrying the data from the
fill message; the array write claims the port afterwards, in parallel. That write is the real
``L2 access'' of a miss, \cref{fig:request-timeline} draws it below the axis, and it lands in no
report column (the \code{ARRAY} lane of the visualizer shows it).

\paragraph{L2-DRAM Bus and DRAM depend on the memory model.} With \code{MainMemoryController}
the DRAM stamps exist and \col{L2-DRAM Bus} is the two memory-bus legs either side of the DRAM
service. MCsim stamps no DRAM events, so there \col{L2-DRAM Bus} becomes ``LLC admit
$\rightarrow$ read leaves on \code{bus[1]}'' and \col{DRAM} becomes the whole \code{bus[1]} round
trip, return leg included. Nothing of the round trip is lost, but the split moves.

\paragraph{A cache-to-cache supply substitutes the responder.} When the LLC never emits because
the line's owner answered on the bus, the last \stamp{L1.SERVICE}/\stamp{L1.EXIT} strictly before
\stamp{RESP\_BUS.EXIT} become the responder milestones, and the responder's \stamp{SERVICE}
replaces \stamp{LLC.ENTER} as the admission anchor --- the owner snoops the request off the bus
and may answer before the LLC admits it, so the LLC's admission is not on the data path. Then
\col{L2 Stall} is the wait for the owner, \col{L2 Access} its service, and \col{Response Bus}
still \stamp{RESP\_BUS.EXIT} $-$ responder \stamp{EXIT}. If no responder stamp exists at all, the
whole admit-to-grant interval is folded into \col{L2 Access} so the columns keep tiling, and the
run's stderr counts the request as \code{noresp}.

\section{Checking a column}

The columns tile to \col{Total} on every row by construction, and the simulator says so: the
\code{EVENT-PATH} line on stderr at the end of a run reports how many requests tiled, how
many did not, and how many had no responder. When a column looks wrong,
\env{OCTOPUS\_EVENT\_DUMP}\code{=N} prints the raw timeline --- every stamp with its role and
cycle --- of the first $N$ requests whose Total exceeds 100 cycles or that crossed the memory
bus; \env{OCTOPUS\_EVENT\_DUMP\_MIN}\code{=<cycles>} raises the threshold to isolate one outlier
row seen in a report. The seven \emph{mechanism tracker} columns that follow \col{Oldest} in the
report (the line's state when the request arrived at the LLC, whether the per-line gate held
it, how many younger responses and array accesses went ahead of it) attribute a stall to its
cause; \cref{ch:monitoring} defines them.

\begin{wherebox}
\begin{itemize}[nosep,leftmargin=1.2em]
\item \file{src/Logger.cpp} --- the stamps, \code{finalizeEvents} and the ``Report: the event timeline is the sole decomposition'' block that computes every column
\item \file{header/Logger.h} --- the \code{Role} and \code{Phase} enums
\item \file{docs/Logger.md} --- the same material as a reference
\item \file{tutorial/00-setup-and-first-run/README.md} --- reading a first report
\end{itemize}
\end{wherebox}
